\section{Introduzione}

\subsection{Contenuto del documento}

La specifica tecnica è un documento che descrive in dettaglio l'architettura, le tecnologie e i pattern utilizzati 
per lo sviluppo del prodotto. Questo documento è destinato a fornire una guida chiara per quella che è la struttura 
e le componenti di esso. \\
Deve essere visto come un riferimento per la progettazione e la codifica del prodotto, in modo da garantire coerenza 
con il \textit{Proof of Concept}.

\subsubsection{Struttura del documento}

Il documento è suddiviso in 4 macro-sezioni:

\begin{itemize}
    \item \textbf{Introduzione}: espone il contenuto e lo scopo del documento, oltre che i Riferimenti Normativi e Informativi;
    \item \textbf{Tecnologie}: descrive le tecnologie utilizzate per la codifica del prodotto, con una breve descrizione di ognuna;
    \item \textbf{\textit{Pattern}}: descrive i \textit{pattern} utilizzati per la progettazione del prodotto, con una breve descrizione di ognuno;
    \item \textbf{Architettura logica}: descrive l'architettura logica del prodotto;
    \item \textbf{Architettura di \textit{deployment}}: descrive in che modo vengono distribuiti i componenti del prodotto;
    \item \textbf{Dettagli progettuali}: descrivono le micro-architetture dei componenti del prodotto, con diagrammi e descrizioni dettagliate.
\end{itemize}

\subsection{Glossario}
All'interno del documento possono essere presenti termini il cui significato può risultare ambiguo o sconosciuto, essi sono 
riportati in italico e alla loro prima occorrenza è associata una \textit{G} a pedice ad indicare la presenza del suddetto termine 
all'interno del \glossario{glossario} che è disponibile a seguente link: \href{https://snakebyteteam.github.io/glossary.html}{Glossario}.

\subsubsection{Riferimenti Normativi}
\begin{itemize}
    \item \textbf{Norme di Progetto 2.0.0}:
          \\ \url{https://snakebyteteam.github.io/3-PB/Documenti_interni/Norme_di_progetto/NdP_v2.0.0.pdf}
          \\(consultato il 5/04/2026).

    \item \textbf{Vimar View4Life Capitolato di Ingegneria del Software Università di Padova 2025 - 2026:}
          \\ \url{https://www.math.unipd.it/~tullio/IS-1/2025/Progetto/C9.pdf}
          \\ (consultato il 5/04/2026).
\end{itemize}

\subsubsection{Riferimenti Informativi}
\begin{itemize}
    \item \textbf{Glossario}:
        \\ \url{https://snakebyteteam.github.io/3-PB/Documenti_interni/Glossario/Glossario_v2.0.0.pdf} 
        \\ (consultato il 5/04/2026).
\end{itemize}